\documentclass[10pt,a4paper]{amsart}
\usepackage{amsmath,amssymb,amsthm,graphicx}
\usepackage{verbatim}
\newenvironment{reference}{\comment}{\endcomment}
%\newenvironment{reference}{}{}
\newenvironment{slogan}{\comment}{\endcomment}
\newenvironment{history}{\comment}{\endcomment}

% For commutative diagrams we use Xy-pic
\usepackage[all]{xy}

% We use 2cell for 2-commutative diagrams.
\xyoption{2cell}
% Theorem environments.
%
\theoremstyle{plain}
\newtheorem{theorem}[subsection]{Theorem}
\newtheorem{proposition}[subsection]{Proposition}
\newtheorem{lemma}[subsection]{Lemma}

\theoremstyle{definition}
\newtheorem{definition}[subsection]{Definition}
\newtheorem{example}[subsection]{Example}
\newtheorem{exercise}[subsection]{Exercise}
\newtheorem{situation}[subsection]{Situation}

\theoremstyle{remark}
\newtheorem{remark}[subsection]{Remark}
\newtheorem{remarks}[subsection]{Remarks}

\numberwithin{equation}{subsection}

% Macros
%
\def\lim{\mathop{\mathrm{lim}}\nolimits}
\def\colim{\mathop{\mathrm{colim}}\nolimits}
\def\Spec{\mathop{\mathrm{Spec}}}
\def\Hom{\mathop{\mathrm{Hom}}\nolimits}
\def\Ext{\mathop{\mathrm{Ext}}\nolimits}
\def\SheafHom{\mathop{\mathcal{H}\!\mathit{om}}\nolimits}
\def\SheafExt{\mathop{\mathcal{E}\!\mathit{xt}}\nolimits}
\def\Sch{\mathit{Sch}}
\def\Mor{\mathop{\mathrm{Mor}}\nolimits}
\def\Ob{\mathop{\mathrm{Ob}}\nolimits}
\def\Sh{\mathop{\mathit{Sh}}\nolimits}
\def\NL{\mathop{N\!L}\nolimits}
\def\CH{\mathop{\mathrm{CH}}\nolimits}
\def\proetale{{pro\text{-}\acute{e}tale}}
\def\etale{{\acute{e}tale}}
\def\QCoh{\mathit{QCoh}}
\def\Ker{\mathop{\mathrm{Ker}}}
\def\Im{\mathop{\mathrm{Im}}}
\def\Coker{\mathop{\mathrm{Coker}}}
\def\Coim{\mathop{\mathrm{Coim}}}

% Boxtimes
%
\DeclareMathSymbol{\boxtimes}{\mathbin}{AMSa}{"02}

%
% Macros for moduli stacks/spaces
%
\def\QCohstack{\mathcal{QC}\!\mathit{oh}}
\def\Cohstack{\mathcal{C}\!\mathit{oh}}
\def\Spacesstack{\mathcal{S}\!\mathit{paces}}
\def\Quotfunctor{\mathrm{Quot}}
\def\Hilbfunctor{\mathrm{Hilb}}
\def\Curvesstack{\mathcal{C}\!\mathit{urves}}
\def\Polarizedstack{\mathcal{P}\!\mathit{olarized}}
\def\Complexesstack{\mathcal{C}\!\mathit{omplexes}}
% \Pic is the operator that assigns to X its picard group, usage \Pic(X)
% \Picardstack_{X/B} denotes the Picard stack of X over B
% \Picardfunctor_{X/B} denotes the Picard functor of X over B
\def\Pic{\mathop{\mathrm{Pic}}\nolimits}
\def\Picardstack{\mathcal{P}\!\mathit{ic}}
\def\Picardfunctor{\mathrm{Pic}}
\def\Deformationcategory{\mathcal{D}\!\mathit{ef}}
\usepackage{cleveref}
\crefname{originallemma}{original lemma}{original lemmas}
\Crefname{originallemma}{Original lemma}{Original lemmas}
\makeatletter
\newlabel{categories-lemma-adjoint-fully-faithful}{{07RB}{source}{}{sourcepointer-1}{}}
\newlabel{categories-lemma-adjoint-fully-faithful@cref}{{[originallemma][0][]07RB}{[1][0][]source}{}{}{}}
\makeatother
\crefname{originalreference}{original reference}{original references}
\Crefname{originalreference}{Original reference}{Original references}
\makeatletter
\newlabel{lemma-derived-restriction-exts}{{0CS5}{source}{}{sourcepointer-2}{}}
\newlabel{lemma-derived-restriction-exts@cref}{{[originalreference][0][]0CS5}{[1][0][]source}{}{}{}}
\makeatother
\crefname{originalreference}{original reference}{original references}
\Crefname{originalreference}{Original reference}{Original references}
\makeatletter
\newlabel{lemma-tensor-hom-adjoint}{{09LT}{source}{}{sourcepointer-3}{}}
\newlabel{lemma-tensor-hom-adjoint@cref}{{[originalreference][0][]09LT}{[1][0][]source}{}{}{}}
\makeatother
\crefname{originalreference}{original reference}{original references}
\Crefname{originalreference}{Original reference}{Original references}
\makeatletter
\newlabel{remark-P-resolution}{{09PB}{source}{}{sourcepointer-4}{}}
\newlabel{remark-P-resolution@cref}{{[originalreference][0][]09PB}{[1][0][]source}{}{}{}}
\makeatother
\crefname{originalreference}{original reference}{original references}
\Crefname{originalreference}{Original reference}{Original references}
\makeatletter
\newlabel{remark-shift-tensor-no-sign}{{0FQN}{source}{}{sourcepointer-5}{}}
\newlabel{remark-shift-tensor-no-sign@cref}{{[originalreference][0][]0FQN}{[1][0][]source}{}{}{}}
\makeatother
\title{Derived tensor with a compact faithful bimodule is fully faithful}
\author{The Stacks Project Authors}
\date{stacks-project:dga.tex:local-snapshot-20261012; GNU FDL 1.2 or later}
\begin{document}
\maketitle
\noindent\textit{Editorial note.} All delivered mathematical text is the author's own, copied from the pinned Stacks Project chapter file sources/stacks-dga/pinned-original-source.tex; no statement or proof step was rewritten, reordered or removed.\par\medskip
\noindent\textit{Editorial note.} The counted proof is the author's complete argument: assumption (2) is the substantive input and is used exactly once, at the shift case; dependencies carried as navigation stubs: Lemma lemma-tensor-hom-adjoint (adjunction), Lemma lemma-derived-restriction-exts (Ext interpretation), Remark remark-P-resolution and Remark remark-shift-tensor-no-sign, all in-source; the adjoint fully-faithful criterion is standard external tooling.\par\medskip
\noindent\textit{Editorial note.} The label lemma-tensor-with-compact-fully-faithful names the counted statement, and the proof environment immediately following it in the source is delivered verbatim as its proof.\par\medskip
\noindent\textit{Editorial note.} The counted statement is carried by the author's own lemma environment in the chapter source.\par\medskip
\noindent\textit{Allowed standard prerequisites (editorial).}\begin{itemize}
\item Tensor-Hom adjunction. Same unit, Lemma lemma-tensor-hom-adjoint (notation and assumptions as stated there). The proof uses that the functor - tensor\_A\textasciicircum{}L N has left adjoint RHom(N, -), given by Lemma lemma-tensor-hom-adjoint.
\item Property (P) resolution induction. Same unit, Remark remark-P-resolution and Remark remark-shift-tensor-no-sign. The reduction to M = A[k] uses the property (P) resolution technique of Remark remark-P-resolution, and the identification A[k] tensor\_A\textasciicircum{}L N = N[k] uses Remark remark-shift-tensor-no-sign.
\end{itemize}\medskip
\ifdefined\hypertarget\hypertarget{sourcepointer-1}{}\fi
\noindent\textit{Source reference.} Original lemma 07RB: sources/stacks-chow/tags.txt, source lines 9861--9861. Prerequisite pointer invoked at proof line(s) 5700; it names a standard fact of a sibling chapter that the counted argument uses, and every step of the counted proof built on it is retained verbatim in the delivered body.\par\medskip
\ifdefined\hypertarget\hypertarget{sourcepointer-2}{}\fi
\noindent\textit{Source reference.} Original reference 0CS5: sources/stacks-dga/pinned-original-source.tex, source lines 5354--5354. Prerequisite pointer invoked at proof line(s) 5711; it names a standard fact of the same chapter that the counted argument uses, and every step of the counted proof built on it is retained verbatim in the delivered body.\par\medskip
\ifdefined\hypertarget\hypertarget{sourcepointer-3}{}\fi
\noindent\textit{Source reference.} Original reference 09LT: sources/stacks-dga/pinned-original-source.tex, source lines 5633--5633. Prerequisite pointer invoked at proof line(s) 5687; it names a standard fact of the same chapter that the counted argument uses, and every step of the counted proof built on it is retained verbatim in the delivered body.\par\medskip
\ifdefined\hypertarget\hypertarget{sourcepointer-4}{}\fi
\noindent\textit{Source reference.} Original reference 09PB: sources/stacks-dga/pinned-original-source.tex, source lines 2953--2953. Prerequisite pointer invoked at proof line(s) 5713; it names a standard fact of the same chapter that the counted argument uses, and every step of the counted proof built on it is retained verbatim in the delivered body.\par\medskip
\ifdefined\hypertarget\hypertarget{sourcepointer-5}{}\fi
\noindent\textit{Source reference.} Original reference 0FQN: sources/stacks-dga/pinned-original-source.tex, source lines 4998--4998. Prerequisite pointer invoked at proof line(s) 5716; it names a standard fact of the same chapter that the counted argument uses, and every step of the counted proof built on it is retained verbatim in the delivered body.\par\medskip
\begin{lemma}
\label{lemma-tensor-with-compact-fully-faithful}
With notation and assumptions as in Lemma \ref{lemma-tensor-hom-adjoint}.
Assume
\begin{enumerate}
\item $N$ defines a compact object of $D(B, \text{d})$, and
\item the map $H^k(A) \to \Hom_{D(B, \text{d})}(N, N[k])$ is an
isomorphism for all $k \in \mathbf{Z}$.
\end{enumerate}
Then the functor $-\otimes_A^\mathbf{L} N$ is fully faithful.
\end{lemma}

\begin{proof}
Our functor has a left adjoint given by
$R\Hom(N, -)$ by Lemma \ref{lemma-tensor-hom-adjoint}.
By Categories, Lemma \ref{categories-lemma-adjoint-fully-faithful}
it suffices to show that for a differential graded $A$-module $M$
the map
$$
M \longrightarrow R\Hom(N, M \otimes_A^\mathbf{L} N)
$$
is an isomorphism in $D(A, \text{d})$. For this it suffices to show that
$$
H^n(M) \longrightarrow
\text{Ext}^n_{D(B, \text{d})}(N, M \otimes_A^\mathbf{L} N)
$$
is an isomorphism, see Lemma \ref{lemma-derived-restriction-exts}.
Since $N$ is a compact object the right hand side commutes
with direct sums. Thus by Remark \ref{remark-P-resolution}
it suffices to prove this map is an isomorphism for $M = A[k]$.
Since $(A[k] \otimes_A^\mathbf{L} N) = N[k]$ by
Remark \ref{remark-shift-tensor-no-sign},
assumption (2) on $N$ is that the result holds for these.
\end{proof}

\end{document}
